\section{Related Work}

\subsection{Deepfake Datasets}
\subsection{Benchmark Datasets}
FaceForensics++ provides manipulated facial videos and a benchmark for evaluating forgery-detection methods \cite{rossler2019faceforensics}. Celeb-DF was introduced as a challenging dataset containing higher-quality synthesized videos, supporting evaluation under more realistic manipulation conditions \cite{li2020celebdf}. Dataset choice and partitioning are consequential: shared identities, source videos, or related frames across training and test partitions can inflate estimates of generalization. Cross-dataset evaluation is therefore useful for measuring domain shift, provided label mapping and preprocessing are documented.

\subsection{CNN-Based Detection}
CNN-based systems learn visual patterns that may distinguish manipulated from authentic media. Xception uses depthwise separable convolutions and has been applied to image classification \cite{chollet2017xception}. Residual networks use identity shortcuts to support deeper CNNs \cite{he2016deep}. MesoNet was designed as a compact network for facial video forgery detection \cite{afchar2018mesonet}. These architectures motivate CNN feature extraction for facial forgery analysis. Comparisons remain meaningful only when datasets, splits, preprocessing, and evaluation units are consistently reported.

\subsection{Positioning of This Study}
\subsection{Attention and Explainability}
The Convolutional Block Attention Module sequentially applies channel and spatial attention to intermediate CNN features, allowing the network to adaptively emphasize informative feature maps and locations \cite{woo2018cbam}. Grad-CAM uses gradients flowing into a convolutional layer to produce a coarse localization map that indicates image regions influential to a class prediction \cite{selvaraju2017gradcam}. In a deepfake detection system, these methods can support attention-guided feature learning and visual inspection of predictions, respectively. A heatmap is an explanation of model sensitivity, however, and should not be interpreted as proof that a highlighted region contains a manipulation artifact.

\subsection{Research Gap and Scope}
Prior work establishes CNN-based detection and attention modules as relevant tools, while the benchmark datasets enable evaluation across manipulation sources. This project focuses on bringing several evaluation dimensions into one proposed application: a controlled comparison with an otherwise equivalent CNN without CBAM, cross-dataset testing using FaceForensics++ and Celeb-DF, robustness tests under compression and reduced resolution, and Grad-CAM visualization. This combination is the scope of the proposed framework, not a claim that each component is unprecedented. Its empirical value remains to be established through completed, consistently controlled experiments.
